% Answers to Chapter 8, from lectures/L08/appendix/c_solutions.md; the self-test from
% lectures/L08/README.md.

\facitavsnitt{c8}{Kapitel 8}{Funktioner}
\begin{facit}

\svar{1.1} \sv{a} Funktion: varje $x$ ger exakt ett $y$.
\sv{b} Ingen funktion: cirkeln ger till exempel både $y = 4$ och $y = -4$ för $x = 0$.
\sv{c} Funktion: alla $x$-värden är olika (att $y = 4$ förekommer två gånger är tillåtet).
\sv{d} Ingen funktion: $x = 1$ ger både $y = 2$ och $y = 5$.

\svar{1.2} \sv{a} $D_f = \mathbb{R} \setminus \{3\}$, $V_f = \mathbb{R} \setminus \{0\}$
\sv{b} $D_g = [-4, \infty)$, $V_g = [0, \infty)$
\sv{c} $D_h = \mathbb{R}$, $V_h = (-\infty, 2]$

\svar{1.3} \sv{a} $k = 1{,}5\,\Omega/{}^{\circ}\text{C}$, $m = 100\,\Omega$
\sv{b} $137{,}5\,\Omega$

\svar{2.1} \sv{a} $f(10) \approx 864$: efter $10$ år finns ungefär $864$ anläggningar.
\sv{b} Efter ungefär $9$ år.

\svar{2.2} \sv{a} En potensfunktion, $P(I) = 50I^2$ (en andragradsfunktion).
\sv{b} $P(0{,}5) = 12{,}5\,\text{W}$, $P(2) = 200\,\text{W}$
\sv{c} $D_P = [0, 3]\,\text{A}$, $V_P = [0, 450]\,\text{W}$

\svar{2.3} \sv{a} Spänningen minskar med $3\,\text{V}$ varje sekund; en konstant förändring ger
ett linjärt samband.
\sv{b} $U(t) = -3t + 12$
\sv{c} $t = 5\,\text{s}$. Nej: spänningen över en urladdande kondensator blir inte negativ utan
planar ut mot $0\,\text{V}$, så modellen gäller bara inom mätintervallet
$0 \leq t \leq 4\,\text{s}$.

\svar{3.1} \sv{a} $7$ \sv{b} $-11$ \sv{c} $2$ \sv{d} $11$ \sv{e} $2$ \sv{f} $-3$

\svar{3.2} \sv{a} $x = 5$ \sv{b} $x = \frac{5}{3} \approx 1{,}67$ \sv{c} $x = 3$ eller $x = -3$
\sv{d} Nej: $x^2 + 2 = 0$ ger $x^2 = -2$, och ingen reell kvadrat är negativ.
\sv{e} $x = 12$

\svar{3.3} \sv{a} $\mathbb{R}$ \sv{b} $\mathbb{R} \setminus \{-2\}$ \sv{c} $[3, \infty)$
\sv{d} $\mathbb{R} \setminus \{-3, 3\}$ \sv{e} $(-\infty, 9]$

\svar{3.4} \sv{a} $[1, \infty)$ \sv{b} $(-\infty, 0]$ \sv{c} $[0, \infty)$ \sv{d} $[-3, 7]$
\sv{e} $[0, \infty)$

\svar{3.5} \sv{a} $f(x) = 3x + 4$ \sv{b} $f(x) = 3x + 2$ \sv{c} $f(x) = -2x + 3$
\sv{d} $f(x) = -3$, en konstant funktion.

\svar{3.6} \sv{a} Lutningen: hur mycket $f(x)$ ändras när $x$ ökar med ett steg.
\sv{b} Skärningspunkten med $y$-axeln, $f(0) = m$.
\sv{c} Avtagande, eftersom $k = -4 < 0$.
\sv{d} $x = 3$
\sv{e} $m = 12\,\text{V}$ är batteriets tomgångsspänning, klämspänningen vid $I = 0$. Lutningen
$k = -2$ är minus den inre resistansen, $R_{\text{i}} = 2\,\Omega$.

\svar{3.7} \sv{a} $f(-3) = 9$, $g(-3) = -27$
\sv{b} En jämn exponent tar bort minustecknet och en udda bevarar det: $(-x)^2 = x^2$ men
$(-x)^3 = -x^3$.
\sv{c} $V_f = [0, \infty)$, $V_g = \mathbb{R}$
\sv{d} $9$ gånger större.

\svar{3.8} \sv{a} $5$, $10$ och $40$ \sv{b} $100$, $50$ och $12{,}5$
\sv{c} $f$ är växande och $g$ avtagande. Basen avgör: $a > 1$ ger en växande funktion,
$0 < a < 1$ en avtagande.
\sv{d} $C = 5$ respektive $C = 100$; $C$ är startvärdet $f(0)$, eftersom $a^0 = 1$.

\svar{3.9} \sv{a} $1{,}05$ \sv{b} $0{,}80$ \sv{c} $f(n) = 2 \cdot 0{,}7^n\,\text{V}$
\sv{d} $0{,}686\,\text{V}$ \sv{e} Efter $2$ steg, då den är $0{,}98\,\text{V}$.

\svar{3.10} \sv{a} $f$: differensen mellan intilliggande värden är konstant, $2$.
\sv{b} $g$: kvoten mellan intilliggande värden är konstant, $2$.
\sv{c} $f(x) = 2x + 3$ \sv{d} $g(x) = 3 \cdot 2^x$ \sv{e} $f(10) = 23$, $g(10) = 3072$

\svar{3.11} \sv{a} $x = 3$ \sv{b} $x = \pm 5$ \sv{c} $x_1 = 2$, $x_2 = 3$
\sv{d} Inga reella nollställen. \sv{e} $x_1 = 0$, $x_2 = 4$

\svar{3.12} \sv{a} $0\,\text{V}$, $6\,\text{V}$ och $10{,}8\,\text{V}$
\sv{b} $V_U = [0;\ 10{,}8]\,\text{V}$
\sv{c} $x = 2\,\text{k}\Omega$
\sv{d} Nej: nämnaren $1 + x$ är alltid större än täljaren $x$, så $U < 12\,\text{V}$. Spänningen
närmar sig $12\,\text{V}$ men når det aldrig.

\svar{3.13} \sv{a} $10\,\text{V}$ \sv{b} $u(2) \approx 3{,}68\,\text{V}$,
$u(4) \approx 1{,}35\,\text{V}$ \sv{c} Avtagande. \sv{d} $V_u = (0;\ 10]\,\text{V}$
\sv{e} Nej: $e^{-t/2} > 0$ för alla $t$, så spänningen närmar sig $0\,\text{V}$ men blir aldrig
exakt noll.

\svar{3.14} \sv{a} Falskt: en kvadrat är aldrig negativ, så $V_f = [0, \infty)$.
\sv{b} Falskt: $x = 0$ ger division med noll, så $D_f = \mathbb{R} \setminus \{0\}$.
\sv{c} Falskt: ett nollställe kräver $f(x) = 0$, men $f(2) = 5$.
\sv{d} Falskt: $x = 0$ ger både $y = 5$ och $y = -5$.
\sv{e} Falskt: startvärdet är $f(0) = 3$; $2$ är förändringsfaktorn.

\svar[Testa dig själv]{T} \sv{1} Till exempel $f(x) = \sqrt{x}$ med $D_f = [0, \infty)$: roten ur
ett negativt tal är inte definierad. (Eller $f(x) = \frac{1}{x}$, där $x = 0$ ger division med
noll.)
\sv{2} $f(2) = 9$, $f(-1) = 6$

\end{facit}
